\section{Разбиение на документы и выделение текста}

За документ принята одна веб-страница --- статья о фильме, о сериале, биография или
авторский материал. Такое разбиение напрашивается само: каждая сущность на обеих
площадках уже живёт по собственному адресу.

Для анализа вручную сохранено 12 страниц: 7 из Википедии и 5 с Кино-Театр.Ру. Брались
документы разных типов и разного объёма --- большие статьи, две заглушки, сериал,
биографии, авторская статья.

Текст выделялся так: берётся содержимое основного контейнера, из него выбрасываются
\texttt{<script>}, \texttt{<style>} и таблицы, снимаются остальные теги, раскрываются
HTML-сущности, схлопываются пробелы.

Способ грубый. Выбрасывая таблицы, мы теряем вместе с версткой и карточку фильма с
полезными полями (год, страна, режиссёр); зато разделы \emph{Примечания} и
\emph{Ссылки} в текст попадают, хотя это библиографический шум. Для оценки корпуса
этого достаточно, при индексации правила придётся уточнять.

\pagebreak
